\documentclass[12pt,a4paper]{article}

% --- Metadata ---
\def\courseshortheader{HỌC MÁY ỨNG DỤNG --- PHÂN LOẠI}
\def\coverbookline{TRÍ TUỆ NHÂN TẠO}
\def\covermaintitle{Phân loại}
\def\coversubtitle{Lý thuyết --- Ngưỡng, metric, ROC/AUC, đa lớp}

% Shared preamble for nhom_5 (requires \courseshortheader before \input)
\usepackage[utf8]{inputenc}
\usepackage[vietnamese]{babel}
\usepackage{amsmath,amssymb,amsthm}
\usepackage{geometry}
\usepackage{enumitem}
\usepackage[most]{tcolorbox}
\usepackage{hyperref}
\usepackage{fancyhdr}
\usepackage{graphicx}
\usepackage{booktabs}
\usepackage{tikz}
\usetikzlibrary{calc,arrows.meta,decorations.pathreplacing,positioning}
\usepackage{eso-pic}
\usepackage{xcolor}

\geometry{a4paper, margin=2cm, bottom=2.5cm}
\hypersetup{colorlinks=true, linkcolor=blue!70!black, urlcolor=blue}
\setlist[itemize]{topsep=4pt, itemsep=2pt}
\setlist[enumerate]{topsep=4pt, itemsep=2pt}

\AddToShipoutPictureFG{%
  \AtPageCenter{%
    \begin{tikzpicture}[remember picture, overlay]
      \node[rotate=45, scale=1, opacity=0.06]{%
        \fontsize{60}{72}\selectfont\bfseries\sffamily Đặng Tấn Quí%
      };
    \end{tikzpicture}%
  }%
}

\pagestyle{fancy}
\fancyhf{}
\fancyhead[L]{\small\textit{\courseshortheader}}
\fancyhead[R]{\small\textit{Đặng Tấn Quí}}
\fancyfoot[C]{\thepage}
\fancyfoot[L]{\footnotesize\textcopyright\ Đặng Tấn Quí}
\fancyfoot[R]{\footnotesize SĐT: 0377\,122\,210}
\renewcommand{\headrulewidth}{0.4pt}
\renewcommand{\footrulewidth}{0.4pt}

\fancypagestyle{plain}{%
  \fancyhf{}
  \fancyfoot[C]{\thepage}
  \fancyfoot[L]{\footnotesize\textcopyright\ Đặng Tấn Quí}
  \fancyfoot[R]{\footnotesize SĐT: 0377\,122\,210}
  \renewcommand{\headrulewidth}{0pt}
  \renewcommand{\footrulewidth}{0.4pt}
}

\newtcolorbox{dongco}[1][]{%
  enhanced, breakable,
  colback=teal!4!white, colframe=teal!60!black,
  fonttitle=\bfseries, title={Động cơ học -- Tại sao cần khái niệm này?}, #1%
}
\newtcolorbox{ynghia}[1][]{%
  enhanced, breakable,
  colback=purple!3!white, colframe=purple!50!black,
  fonttitle=\bfseries, title={Ý nghĩa \& Trực giác}, #1%
}
\newtcolorbox{ungdung}[1][]{%
  enhanced, breakable,
  colback=olive!5!white, colframe=olive!60!black,
  fonttitle=\bfseries, title={Ứng dụng thực tế}, #1%
}
\newtcolorbox{dinhnghia}[1][]{%
  enhanced, breakable,
  colback=blue!3!white, colframe=blue!60!black,
  fonttitle=\bfseries, title={Định nghĩa}, #1%
}
\newtcolorbox{dinhly}[1][]{%
  enhanced, breakable,
  colback=green!3!white, colframe=green!50!black,
  fonttitle=\bfseries, title={Định lý}, #1%
}
\newtcolorbox{tinhchat}[1][]{%
  enhanced, breakable,
  colback=yellow!5!white, colframe=orange!50!black,
  fonttitle=\bfseries, title={Tính chất}, #1%
}
\newtcolorbox{phuongphap}[1][]{%
  enhanced, breakable,
  colback=gray!5!white, colframe=gray!50!black,
  fonttitle=\bfseries, title={Phương pháp}, #1%
}
\newtcolorbox{luuy}[1][]{%
  enhanced, breakable,
  colback=red!3!white, colframe=red!50!black,
  fonttitle=\bfseries, title={Lưu ý}, #1%
}
\newtcolorbox{congthuc}[1][]{%
  enhanced, breakable,
  colback=cyan!3!white, colframe=cyan!50!black,
  fonttitle=\bfseries, title={Công thức}, #1%
}
\newtcolorbox{vidu}[1][]{%
  enhanced, breakable,
  colback=violet!3!white, colframe=violet!50!black,
  fonttitle=\bfseries, title={Ví dụ}, #1%
}
\newtcolorbox{phanvidu}[1][]{%
  enhanced, breakable,
  colback=red!2!white, colframe=red!40!black,
  fonttitle=\bfseries, title={Phản ví dụ}, #1%
}
\newtcolorbox{tongket}[1][]{%
  enhanced, breakable,
  colback=blue!4!white, colframe=blue!70!black,
  fonttitle=\bfseries, title={Tổng kết chương}, #1%
}


\usepackage{tabularx}
\usepackage{array}
\usepackage{float}
\usepackage{amsmath}

\graphicspath{{images/classification/}}

\newcommand{\mlimg}[2][0.88]{%
  \def\mlimgfile{images/classification/#2}%
  \IfFileExists{\mlimgfile}{%
    \begin{center}\includegraphics[width=#1\linewidth]{#2}\end{center}%
  }{%
    \begin{luuy}[title={Thiếu hình minh họa}]%
      Không tìm thấy \expandafter\path\expandafter{\mlimgfile}. Chạy \path{download_mlcc_images.sh} trong \path{_template/}.%
    \end{luuy}%
  }%
}
\newcommand{\clfimg}[2][0.46\linewidth]{%
  \def\mlimgfile{images/classification/#2}%
  \IfFileExists{\mlimgfile}{%
    \includegraphics[width=#1]{#2}%
  }{%
    \fbox{\begin{minipage}[c][2.2cm][c]{#1}\centering\footnotesize Thiếu hình\end{minipage}}%
  }%
}

\begin{document}

\thispagestyle{empty}
\begin{tikzpicture}[remember picture, overlay]
  \fill[blue!80!black] (current page.south west) rectangle (current page.north east);
  \fill[blue!60!cyan, opacity=0.8]
    (current page.south west) -- (current page.north west) --
    ([xshift=-3cm]current page.north east) -- (current page.south east) -- cycle;
  \foreach \r/\o in {6/0.05, 4.5/0.08, 3/0.12} {
    \fill[white, opacity=\o] ([xshift=2cm, yshift=-2cm]current page.north east) circle (\r cm);
  }
  \draw[white, line width=2pt] ([yshift=-5cm]current page.north west) ++(2cm,0) -- ++(17cm,0);
  \draw[white, line width=2pt] ([yshift=-16cm]current page.north west) ++(2cm,0) -- ++(17cm,0);
  \node[anchor=west, white, font=\fontsize{28}{34}\bfseries\selectfont]
    at ([xshift=3cm, yshift=-7.5cm]current page.north west) {\coverbookline};
  \node[anchor=west, white, font=\fontsize{20}{26}\selectfont]
    at ([xshift=3cm, yshift=-9cm]current page.north west) {\covermaintitle};
  \node[anchor=west, white, font=\fontsize{16}{22}\selectfont]
    at ([xshift=3cm, yshift=-10.2cm]current page.north west) {\coversubtitle};
  \node[anchor=west, white, font=\Large\bfseries]
    at ([xshift=3cm, yshift=-13cm]current page.north west) {Biên soạn:};
  \node[anchor=west, yellow!80!white, font=\fontsize{24}{30}\bfseries\selectfont]
    at ([xshift=3cm, yshift=-14.5cm]current page.north west) {ĐẶNG TẤN QUÍ};
  \node[anchor=west, white!90, font=\large]
    at ([xshift=3cm, yshift=-18cm]current page.north west) {SĐT: 0377\,122\,210};
  \node[anchor=south, white!80, font=\small]
    at ([yshift=1.5cm]current page.south) {Tài liệu có bản quyền --- Cấm sao chép khi chưa được sự đồng ý của tác giả};
  \node[anchor=south east, white, font=\Large\bfseries]
    at ([xshift=-2cm, yshift=3cm]current page.south east) {\the\year};
\end{tikzpicture}

\newpage
\tableofcontents
\newpage

\section*{Lời nói đầu}
\addcontentsline{toc}{section}{Lời nói đầu}

\begin{ynghia}[title={Nguồn và cách đọc}]
Tài liệu biên soạn và dịch từ \emph{Classification} (Google ML Crash Course)\footnote{\url{https://developers.google.com/machine-learning/crash-course/classification}}.
\end{ynghia}

\begin{luuy}[title={Điều kiện tiên quyết}]
Cần nắm \href{https://developers.google.com/machine-learning/intro-to-ml}{Introduction to ML}, \href{https://developers.google.com/machine-learning/crash-course/linear-regression}{Linear regression} và \href{https://developers.google.com/machine-learning/crash-course/logistic-regression}{Logistic regression}.
\end{luuy}

\begin{phuongphap}[title={Mục tiêu học tập}]
\begin{itemize}
  \item Xác định ngưỡng phù hợp cho phân loại nhị phân.
  \item Tính và chọn metric đánh giá mô hình.
  \item Đọc và diễn giải ROC và AUC.
\end{itemize}
\end{phuongphap}

\newpage

%==============================================================================
\section{Giới thiệu phân loại}
%==============================================================================

\begin{dongco}[title={Từ 75\% spam sang lớp ``spam''}]
Logistic cho $P(\text{spam})=0{,}75$. Ứng dụng lọc thư cần \textbf{lớp} spam/không spam --- đó là \textbf{phân loại}.
\end{dongco}

\begin{dinhnghia}[title={Phân loại}]
\href{https://developers.google.com/machine-learning/glossary\#classification-model}{Phân loại (classification)}: dự đoán mẫu thuộc \href{https://developers.google.com/machine-learning/glossary\#class}{lớp} nào.
\begin{itemize}
  \item \textbf{Nhị phân:} hai lớp.
  \item \textbf{Đa lớp:} hơn hai lớp (giới thiệu cuối module).
\end{itemize}
\end{dinhnghia}

\begin{phuongphap}[title={Luồng module}]
Xác suất (logistic) $\xrightarrow{\text{ngưỡng}}$ lớp $\xrightarrow{\text{ma trận nhầm lẫn}}$ metric $\xrightarrow{\text{mọi ngưỡng}}$ ROC/AUC.
\end{phuongphap}

\newpage

%==============================================================================
\section{Ngưỡng và ma trận nhầm lẫn}
%==============================================================================

\begin{dongco}[title={Từ 0,75 sang thư mục spam}]
Mô hình spam cho xác suất 0--1. Triển khai cần chuyển \texttt{0.75} thành ``spam'' hoặc ``không spam'' bằng \textbf{ngưỡng}.
\end{dongco}

\begin{dinhnghia}[title={Ngưỡng phân loại}]
\href{https://developers.google.com/machine-learning/glossary\#classification-threshold}{Ngưỡng} $t$: $y' \geq t$ $\Rightarrow$ \href{https://developers.google.com/machine-learning/glossary\#positive_class}{lớp dương} (spam); $y' < t$ $\Rightarrow$ \href{https://developers.google.com/machine-learning/glossary\#negative_class}{lớp âm}. Với thư viện Keras: $y'=t$ xếp vào \textbf{âm}.
\end{dinhnghia}

\begin{vidu}[title={0,51 vs 0,99 với ngưỡng khác nhau}]
Email A: $y'=0{,}99$; email B: $y'=0{,}51$. Ngưỡng $0{,}5$: cả hai là spam. Ngưỡng $0{,}95$: chỉ A là spam.
\end{vidu}

\begin{luuy}[title={0,5 không luôn tối ưu}]
Tránh mặc định 0,5 khi: lớp mất cân bằng (0,01\% spam); chi phí FP/FN không đối xứng; phân loại nhầm thư hợp lệ đắt hơn phân loại spam lọt vào hộp thư đến.
\end{luuy}

\subsection{Ma trận nhầm lẫn}

\begin{dinhnghia}[title={Confusion matrix}]
  Điểm số xác suất không phải là thực tế, cũng không phải là sự thật khách quan (\href{https://developers.google.com/machine-learning/glossary\#ground_truth}{ground truth}). Bốn ô (cột = thực tế, hàng = dự đoán):
  \begin{itemize}
    \item Probability score (Điểm số xác suất): Con số thể hiện khả năng xảy ra của một sự kiện (thường là kết quả dự đoán của thuật toán, mô hình AI).
    \item Reality (Thực tế): Những gì thực sự xảy ra hoặc tồn tại trong đời thực.
    \item Ground truth (Sự thật khách quan/Dữ liệu thực tế): Kết quả chính xác, đúng đắn đã được kiểm chứng ngoài đời thực.
  \end{itemize}
\end{dinhnghia}

\begin{congthuc}[title={Ma trận nhầm lẫn (spam)}]
\begin{tabularx}{\textwidth}{@{}l >{\raggedright\arraybackslash}X >{\raggedright\arraybackslash}X@{}}
 & \textbf{Thực tế dương} & \textbf{Thực tế âm} \\
\midrule
\textbf{Dự đoán dương} &
\textbf{TP (True positive) (Dương tính thật)} --- spam đúng vào thư rác &
\textbf{FP (False positive) (Dương tính giả)} --- thư hợp lệ vào spam \\
\textbf{Dự đoán âm} &
\textbf{FN (False negative) (Âm tính giả)} --- spam lọt hộp thư &
\textbf{TN (True negative) (Âm tính thật)} --- thư hợp lệ đúng hộp thư \\
\end{tabularx}
\end{congthuc}

\begin{tinhchat}[title={Tổng theo hàng và cột}]
\begin{itemize}
  \item \textbf{Hàng:} mọi dự đoán dương ($TP+FP$), mọi dự đoán âm ($FN+TN$).
  \item \textbf{Cột:} mọi thực tế dương ($TP+FN$), mọi thực tế âm ($FP+TN$).
\end{itemize}
\end{tinhchat}

\begin{vidu}[title={Dữ liệu mất cân bằng}]
Tổng thực tế dương không gần bằng tổng âm $\Rightarrow$ \href{https://developers.google.com/machine-learning/glossary\#class_imbalanced_data_set}{mất cân bằng lớp}.
\end{vidu}

\subsection{Ảnh hưởng của ngưỡng}

\begin{ungdung}[title={Video ngưỡng}]
\url{https://www.youtube.com/watch?v=QM0sYbEQSkM}
\end{ungdung}

\begin{luuy}[title={Widget MLCC (Separated / Unseparated / Imbalanced)}]
Kéo ngưỡng từ 0,1 $\rightarrow$ 0,9 để quan sát TP, FP, FN, TN thay đổi.

Liên kết: \href{https://developers.google.com/machine-learning/crash-course/classification/thresholding}{Widget MLCC (Separated / Unseparated / Imbalanced)}
\end{luuy}

\begin{luuy}[title={Ngưỡng}]
Site hợp lệ bị gán malware $\Rightarrow$ \textbf{FP} (âm bị gán dương - dương tính giả).

Tăng ngưỡng $\Rightarrow$ \textbf{TP và FP giảm; còn FN và TN tăng}.
\end{luuy}

\newpage

%==============================================================================
\section{Accuracy, recall, precision và metric liên quan}
%==============================================================================

\begin{dongco}[title={Một con số không đủ}]
Metric phụ thuộc chi phí sai, mất cân bằng lớp và \textbf{ngưỡng} đang dùng (mỗi metric tính tại một ngưỡng cố định).
\begin{itemize}
    \item Accuracy: Độ chính xác
    \item Recall: Độ nhạy / độ thu hồi / tỷ lệ dự báo đúng
    \item Precision: Độ chuẩn xác / độ chuẩn xác dương tính
    \item Metric: Đo lường / chỉ số đánh giá
  \end{itemize}
\end{dongco}

\begin{congthuc}[title={Các công thức chính}]
\[
  \text{Accuracy}=\frac{TP+TN}{TP+TN+FP+FN},\quad
  \text{Recall}=\frac{TP}{TP+FN},\quad
  \text{FPR}=\frac{FP}{FP+TN}
\]
\[
  \text{Precision}=\frac{TP}{TP+FP},\quad
  F_1=\frac{2\cdot\text{precision}\cdot\text{recall}}{\text{precision}+\text{recall}}=\frac{2TP}{2TP+FP+FN}
\]
\end{congthuc}

\begin{ynghia}[title={Diễn giải (ví dụ spam)}]
\begin{itemize}
  \item \textbf{Accuracy:} tỷ lệ email phân loại đúng. Cho biết tổng tỷ lệ dự đoán đúng của toàn bộ mô hình trên mọi trường hợp. Một mô hình hoàn hảo sẽ không có kết quả dương tính giả và không có kết quả âm tính giả, do đó, độ chính xác là 1.0.
  \item \textbf{Recall (TPR - True Positive Rate):} trong spam thật, bắt được bao nhiêu? (\emph{probability of detection - xác suất phát hiện}). Trả lời câu hỏi: "Trong tổng số trường hợp thực tế cần tìm, mô hình đã quét được bao nhiêu phần trăm?". Một mô hình hoàn hảo giả định sẽ không có âm tính giả và do đó có độ thu hồi (TPR) là 1.0.
  \item \textbf{FPR (False Positive Rate):} tỷ lệ thư hợp lệ bị gán nhầm spam. (\emph{probability of false alarm - xác suất báo động giả}). Trả lời câu hỏi: "Trong số những lần mô hình báo động (dự đoán dương tính), có bao nhiêu lần là cảnh báo giả?". Một mô hình hoàn hảo sẽ không có kết quả dương tính giả và do đó có FPR là 0.0.
  \item \textbf{Precision:} trong email gán spam, bao nhiêu đúng là spam? Trả lời câu hỏi: "Mô hình dự đoán đúng được bao nhiêu trong số những lần nó cảnh báo?". Một mô hình hoàn hảo giả định sẽ không có dương tính giả và do đó có độ chính xác là 1.0.
  \item \textbf{F1:} "Mô hình hoạt động cân bằng và ổn định như thế nào?" Trả lời câu hỏi: "Mô hình đang hoạt động ở mức độ nào? Độ nhạy và độ chính xác có đồng đều không?"
\end{itemize}
\end{ynghia}

\begin{phanvidu}[title={Trên tập mất cân bằng}]
  Tuy nhiên, khi tập dữ liệu không cân bằng hoặc khi một loại lỗi (FN hoặc FP) tốn kém hơn loại còn lại (trường hợp này xảy ra trong hầu hết các ứng dụng thực tế), bạn nên tối ưu hoá cho một trong các chỉ số khác. Đối với các tập dữ liệu có độ mất cân bằng cao, trong đó một lớp xuất hiện rất hiếm, chẳng hạn như 1\%, một mô hình dự đoán giá trị âm 100\% sẽ đạt điểm 99\% về độ chính xác, mặc dù không có tác dụng.

  Trong đó số lượng kết quả dương tính thực tế rất thấp, độ thu hồi là một chỉ số có ý nghĩa hơn độ chính xác vì độ thu hồi đo lường khả năng của mô hình trong việc xác định chính xác tất cả các trường hợp dương tính. Đối với các ứng dụng như dự đoán bệnh, việc xác định chính xác các trường hợp dương tính là rất quan trọng. Thông thường, kết quả âm tính giả sẽ có hậu quả nghiêm trọng hơn so với kết quả dương tính giả.

  FPR thường là một chỉ số giàu thông tin hơn độ chính xác. Tuy nhiên, nếu số lượng âm tính thực tế rất thấp, thì FPR có thể không phải là lựa chọn lý tưởng do tính biến động của nó. Ví dụ: nếu chỉ có 4 giá trị âm thực tế trong một tập dữ liệu, thì một lần phân loại sai sẽ dẫn đến FPR là 25\%, trong khi lần phân loại sai thứ hai sẽ khiến FPR tăng lên 50\%. Trong những trường hợp như thế này, precision có thể là một chỉ số ổn định hơn để đánh giá tác động của kết quả dương tính giả.

  Trong đó số lượng kết quả dương tính thực tế rất thấp, chẳng hạn như tổng cộng 1-2 ví dụ, precision sẽ ít có ý nghĩa và ít hữu ích hơn.
\end{phanvidu}

\begin{tinhchat}[title={Precision vs recall}]
  Độ chính xác cải thiện khi số lượng kết quả dương tính giả giảm, trong khi độ thu hồi cải thiện khi số lượng kết quả âm tính giả giảm. Nhưng như đã thấy trong phần trước, việc tăng ngưỡng phân loại có xu hướng làm giảm số lượng kết quả dương tính giả và tăng số lượng kết quả âm tính giả, trong khi việc giảm ngưỡng sẽ có tác dụng ngược lại. Do đó, độ chính xác và khả năng thu hồi thường có mối quan hệ nghịch đảo, trong đó việc cải thiện một trong hai chỉ số này sẽ làm giảm chỉ số còn lại.
\end{tinhchat}

\begin{luuy}[title={NaN}]
$TP=FP=0$ $\Rightarrow$ precision không xác định. Mô hình không bao giờ dự đoán dương có thể cho NaN.
\end{luuy}

\begin{congthuc}[title={Gợi ý chọn metric}]
\begin{tabularx}{\textwidth}{@{}l X@{}}
\textbf{Metric} & \textbf{Khi nào ưu tiên} \\
\midrule
Accuracy & Tập cân bằng; chỉ dùng kèm metric khác. \\
Recall & FN đắt (bỏ sót bệnh, spam lọt). \\
FPR & FP đắt (báo động giả). \\
Precision & Cần dự đoán dương chính xác. \\
$F_1$ & Cân bằng precision/recall; tập mất cân bằng. \\
\end{tabularx}
\end{congthuc}

\begin{vidu}[title={Tính recall và precision}]
$TP=5$, $TN=6$, $FP=3$, $FN=2$:

Recall $=5/7\approx 0{,}714$,

Precision $=5/8=0{,}625$.

$TP=3$, $FP=2$, $FN=1$:

Precision $=3/5=0{,}6$.
\end{vidu}

\begin{luuy}[title={Metric}]
$TP=5$, $FN=2$ $\Rightarrow$ Recall $=5/7$ (không chia cho tổng mọi phân loại đúng).

$TP=3$, $FP=2$ $\Rightarrow$ Precision $=3/5$ (chia cho mọi dự đoán dương).

Bẫy côn trùng xâm lấn: FN đắt, FP rẻ $\Rightarrow$ tối đa \textbf{Recall}.
\end{luuy}

\newpage

%==============================================================================
\section{ROC và AUC}
%==============================================================================

\begin{dongco}[title={Đánh giá mọi ngưỡng}]
Accuracy/recall tại \textbf{một} ngưỡng. ROC quét ngưỡng --- so sánh mô hình toàn cục.
\end{dongco}

\begin{dinhnghia}[title={Đường cong ROC}]
Với mỗi ngưỡng: tính TPR (recall) và FPR; vẽ TPR ($y$) theo FPR ($x$). Điểm $(0,1)$: hoàn hảo (FPR=0, TPR=1).
\end{dinhnghia}

\mlimg[0.55]{auc_1-0.png}
\begin{center}\small\textbf{Hình 1.} ROC mô hình hoàn hảo (AUC = 1).\end{center}

\begin{dinhnghia}[title={AUC}]
\href{https://developers.google.com/machine-learning/glossary\#AUC}{AUC}: xác suất mô hình xếp ngẫu nhiên một mẫu dương \textbf{cao hơn} một mẫu âm (không phụ thuộc ngưỡng).
\end{dinhnghia}

\mlimg[0.72]{no_slider_data.png}
\begin{center}\small\textbf{Hình 2.} AUC = $P$(vuông bên phải hình tròn).\end{center}

\mlimg[0.55]{auc_0-5.png}
\begin{center}\small\textbf{Hình 3.} Đoán ngẫu nhiên: đường chéo, AUC = 0,5.\end{center}

\begin{ynghia}[title={Precision-Recall curve (tùy chọn)}]
Tập mất cân bằng nặng: PRC (precision vs recall) đôi khi so sánh tốt hơn ROC.
\end{ynghia}

\mlimg[0.65]{prauc.png}
\begin{center}\small\textbf{Hình 4.} Ví dụ đường precision-recall.\end{center}

\begin{figure}[H]
\centering
\clfimg{auc_0-65.png}\hfill
\clfimg{auc_0-93.png}
\begin{center}\small\textbf{Hình 5.} AUC = 0,65 (trái) vs 0,93 (phải).\end{center}
\end{figure}

\mlimg[0.72]{auc_abc.png}
\begin{center}\small\textbf{Hình 6.} Điểm A, B, C --- chọn ngưỡng theo chi phí FP/FN.\end{center}

\begin{phuongphap}[title={Chọn điểm trên ROC}]
\begin{itemize}
  \item \textbf{A:} FP đắt $\Rightarrow$ FPR thấp, chấp nhận TPR thấp hơn.
  \item \textbf{C:} FN đắt $\Rightarrow$ tối đa TPR.
  \item \textbf{B:} chi phí tương đương $\Rightarrow$ cân bằng.
\end{itemize}
\end{phuongphap}

\begin{luuy}[title={AUC nhỏ hơn 0,5}]
Tệ hơn ngẫu nhiên --- có thể \textbf{đảo nhãn} 0$\leftrightarrow$1 mà không cần huấn luyện lại.
\end{luuy}

\begin{figure}[H]
\centering
\clfimg[0.22\linewidth]{auc_0-77.png}
\clfimg[0.22\linewidth]{auc_0-508.png}
\clfimg[0.22\linewidth]{auc_0-623.png}
\clfimg[0.22\linewidth]{auc_0-31.png}
\begin{center}\small\textbf{Hình 7.} So sánh AUC: 0,77 tốt nhất; 0,31 tệ hơn ngẫu nhiên.\end{center}
\end{figure}

\begin{luuy}[title={ROC}]
AUC cao nhất: ($\approx 0{,}77$).

Tệ hơn ngẫu nhiên: (AUC nhỏ hơn 0,5).

Cải thiện không huấn luyện lại: \textbf{đảo nhãn dự đoán}.

Thư quan trọng không được vào spam $\Rightarrow$ \textbf{điểm A} (FPR thấp).
\end{luuy}

\newpage

%==============================================================================
\section{Độ lệch dự đoán (prediction bias)}
%==============================================================================

\begin{dongco}[title={Kiểm tra nhanh sau huấn luyện}]
Trung bình xác suất dự đoán có khớp tỷ lệ nhãn thật? Lệch sớm $\Rightarrow$ lỗi dữ liệu, pipeline hoặc chuẩn hóa quá mạnh.
\end{dongco}

\begin{congthuc}[title={Prediction bias}]
\[
  \text{bias} = \bar{y}' - \bar{y}
\]
$\bar{y}$: trung bình nhãn ground truth; $\bar{y}'$: trung bình dự đoán. Ví dụ: 5\% spam $\Rightarrow$ kỳ vọng $\bar{y}' \approx 0{,}05$. Mô hình luôn ``50\% spam'' $\Rightarrow$ lệch nghiêm trọng.
\end{congthuc}

\begin{tinhchat}[title={Nguyên nhân có thể}]
\begin{itemize}
  \item Dữ liệu / lấy mẫu thiên lệch.
  \item Regularization quá mạnh.
  \item Lỗi pipeline huấn luyện.
  \item Đặc trưng không đủ cho tác vụ.
\end{itemize}
\end{tinhchat}

\newpage

%==============================================================================
\section{Phân loại đa lớp}
%==============================================================================

\begin{dongco}[title={Hơn hai lớp}]
Nhận dạng chữ số 0--9; phân loại nhiều loài --- không chỉ spam/không spam.
\end{dongco}

\begin{phuongphap}[title={One-vs-rest}]
Mỗi mẫu một lớp: chia thành nhiều bài toán nhị phân. Ba lớp A, B, C:
\begin{enumerate}
  \item (A+B) vs C.
  \item Trong (A+B): A vs B.
\end{enumerate}
\end{phuongphap}

\begin{vidu}[title={Chữ số viết tay}]
Ảnh một chữ số $\Rightarrow$ dự đoán 0--9.
\end{vidu}

\begin{luuy}[title={Đa nhãn}]
Một mẫu thuộc \textbf{nhiều lớp cùng lúc} $\Rightarrow$ \textbf{đa nhãn} (multi-label), khác đa lớp.
\end{luuy}

\begin{tongket}[title={Tổng kết chương --- Phân loại}]
\begin{itemize}
  \item Xác suất + \textbf{ngưỡng} $\to$ lớp; \textbf{confusion matrix} $\to$ metric.
  \item Chọn metric theo chi phí FN/FP; \textbf{ROC/AUC} cho mọi ngưỡng.
  \item Kiểm tra \textbf{prediction bias}; đa lớp = one-vs-rest hoặc mô hình chuyên sâu.
\end{itemize}
\end{tongket}

\end{document}
